\section{Related Work}
\label{sec:related}

\subsection{Blockchain Consensus and the Bitcoin Backbone}

Garay, Kiayias, and Leonardos~\cite{GKL20} formalise the core of Bitcoin as the \emph{Bitcoin backbone} protocol, operating in the random oracle model under a synchronous network with bounded message delay. Security is characterised by three properties: Common Prefix, Chain Quality, and Chain Growth, proved by bounding per-round random variables via Chernoff concentration within a \emph{typical execution} framework. The honest majority condition requires that the adversary's share of hash power remains below a threshold parameterised by the block production rate~$f$. This framework has become the standard reference for blockchain security proofs, with extensions to variable difficulty~\cite{GKL17}, partial synchrony, and proof-of-stake~\cite{KRDO17}. Our work builds directly on it: we show that PoML's block-production mechanism maps to a $q_{\mathsf{ML}}$-bounded random oracle model, enabling a direct invocation of the backbone theorems with adjusted parameters.

\subsection{Practical Proof-of-Useful-Work}

Recent research has largely focused on replacing energy-intensive hashing with useful machine learning tasks. Lan et al.\ propose Proof of Learning (PoLe)~\cite{pole}, where consensus nodes train neural networks on public data, using a Secure Mapping Layer (SML) to prevent cheating and link blocks. While this redirects energy toward optimisation, the SML adds computational complexity and the protocol is strictly limited to training rather than inference. Similarly, Lihu et al.\ introduce a PoUW protocol~\cite{pouw} where miners derive nonces from training byproducts such as gradient updates, though verification requires re-running training iterations. Krzywiecki and Wechta advance this with Proof of Training (PoT)~\cite{pot}, which uses cryptographic training secrets to verify delegated tasks, but faces challenges around dataset privacy and public availability.

The most rigorous PoUW analyses appear in Ofelimos~\cite{FKPR22} and D-PoDL~\cite{SLT23}. Fitzi et al.\ prove all three backbone properties for a protocol in which miners perform Doubly Parallel Local Search on combinatorial optimisation problems, using a hash--work--hash structure with SNARG-based verification. Security follows cleanly because DPLS is designed to produce memoryless Bernoulli trials by construction. Su, Larangeira, and Tanaka apply the same backbone machinery to deep learning training in D-PoDL, introducing a model-referencing mechanism for collaborative training and proving chain growth, chain quality, and common prefix under an empirical distributional assumption on training accuracy. Both protocols incur non-useful pre-hash overhead and depend on external parties to supply problem instances or training tasks.

Moving beyond standalone training, Sokhankhosh and Rouhani propose Proof-of-Collaborative-Learning (PoCL)~\cite{pocl}, a federated learning consensus where miners evaluate peers' predictions and select winners via voting, trading cryptographic verification for distributed fairness. The closest architectural alignment to our proposal is Oleksak et al.~\cite{zksnarkmarket}, who design a consensus-layer marketplace for outsourcing zk-SNARK generation with unforgeability and freshness guarantees. While they treat proof generation itself as the commodity, PoML integrates ZKPs specifically to attest to the correctness of ML inference, and extends this foundation to a heterogeneous setting with valuation mechanisms for ranking diverse models---a feature absent in homogeneous proof-generation markets. Unlike all training-centric protocols above, our Proof-of-ML-Computation targets inference, eliminating multi-round commitment, distributional assumptions, and external task publishers, as users submit queries directly to a natural on-chain marketplace.

\subsection{Zero-Knowledge Proofs for ML Inference}

Recent advances in Zero-Knowledge Machine Learning (ZKML) have shifted from specialised protocols for small models to robust compilers capable of verifying large language models. Chen et al.\ introduced ZKML~\cite{chen2024zkml}, an optimising compiler transforming TensorFlow models into Halo2 circuits, improving proving speeds by up to $24\times$ over prior work, though it remains constrained by memory usage and latency at scale. To address scalability, the same group developed ZKTorch~\cite{chen2025zktorch}, which uses a parallelised Mira accumulation scheme to verify all MLPerf Edge Inference models with substantially smaller proof sizes. Similarly, Xie et al.\ propose zkPyTorch~\cite{xie2025zkpytorch}, integrating PyTorch with the Expander engine through ZKP-friendly quantisation, enabling verification of models like Llama-3 without deep cryptographic expertise.

Addressing Transformer-specific bottlenecks, Sun et al.\ present zkLLM~\cite{sun2024zkllm}, which introduces specialised protocols for non-arithmetic operations such as Softmax and attention, enabling verification of 13-billion parameter models in under 15 minutes. As surveyed by Peng et al.~\cite{peng2025survey}, efficiency around floating-point quantisation and non-linear functions remains the primary implementation challenge across the field. These technical strides in efficient verifiable inference are foundational to PoML, providing the cryptographic primitives needed to transform model execution into a verifiable consensus metric for a practical decentralised marketplace.
